% =====================================================================================================
\section{Supplement: further results}\label{sec:supp}

Each result below comes with one line on why it matters; the status notes give their standing.

\subsection{The posterior}\label{sec:supp-post}

\subsubsection{Fixed design}\label{sec:supp-fixed}

\why{The same conclusions hold for a deterministic design, under conditions on the observed sample that a user can
check; it is the route for bounded data such as genotypes.}

The theorem uses four assumptions beyond those of Theorem~\ref{thm:post-random}'s setting.

\begin{assumption}{BC}{Bounded coefficients}
$\max_{j}\abs{\theta^*_j}\le\bar\theta$ for a fixed $\bar\theta<\infty$. \hfill [\SA{post}]
\end{assumption}

\begin{assumption}{EF}{Eigenvalue floor on sets of bounded size}
$\lmin(\Rhat_S)\ge\varphi_0^2$ for every model $S$ with $\abs{S}\le2(A+1)\sz+2$, for a fixed $\varphi_0\in(0,1]$. \hfill [\SA{post}]
\end{assumption}

\begin{assumption}{NC}{No near-exact collinearity among sets of logarithmic size}
$\lmin(G_S)\ge g_0$ for every model $S$ with $\abs{S}\le3(A+1)(\sz+1)\log n$, for a fixed $g_0>0$; on the correlation
scale, a floor $g_0/(n-1)$ that vanishes as $n$ grows. \hfill [\SA{post}]
\end{assumption}

\begin{assumption}{BM}{Beta-min}
Either (a) $\tmin\ge C_\beta\,\sigs\rn$ with a fixed $C_\beta>0$ such that $C_\beta^2\varphi_0^2>c_0(\sz,A)$, an explicit
constant (Appendix~\ref{app:constants}); or (b) $\tmin\sqrt{n/\log n}\to\infty$. \hfill [\SA{post}]
\end{assumption}

\begin{theorem}[Posterior concentration, fixed design]\label{thm:post-fixed}
Let the design be fixed and assume \as{GL}, \as{DIM}, \as{SP}, \as{BC}, \as{EF}, \as{NC} and \as{BM}. Then, along every
sequence that satisfies these assumptions,
\begin{equation}\label{eq:post-op}
  \max_{j\in\set{p}}\big|\PIP_j-\ind{j\in\Sz}\big|\ \le\ 1-\pis{\Sz}\ =\ \OP\Big(\sqrt{\frac{\log n}{n}}\Big),\qquad
  \P\big(\HPM=\MPM=\Sz\big)\to1 .
\end{equation}
Moreover, there are $n_0$ and $C_{\mathrm E}$, depending only on $(\sz,A,\sigs^2,\bar\theta,\varphi_0,g_0,C_\beta)$, such that for
every $n\ge n_0$ and every design and $\tstar$ that satisfy the assumptions at that $n$,
\begin{equation}\label{eq:post-exp}
  \E\big[1-\pis{\Sz}\big]\ \le\ C_{\mathrm E}\sqrt{\frac{\log n}{n}},\qquad \P\big(\HPM=\MPM=\Sz\big)\ \ge\ 1-4C_{\mathrm E}\sqrt{\frac{\log n}{n}} .
\end{equation}
The expectation is over the noise, with the design fixed. Under form (b) of \as{BM}, $n_0$ also depends on how fast
$\tmin\sqrt{n/\log n}$ diverges.
\end{theorem}


\status{The in-probability conclusions~\eqref{eq:post-op}: \textbf{proved and independently checked}. The expectation
bound~\eqref{eq:post-exp}: \textbf{proved and independently checked}: its last unchecked input (a proposition giving a lower
bound on the posterior miss) was recorded as independently checked on 2026-11-17, after the freeze. The registry string
still reads ``conditional on'' one input, a lemma on random Gaussian designs that is itself proved and independently
checked (Appendix~\ref{app:status-open}, housekeeping). Registry: Appendix~\ref{app:status}, row
1.1 (the ledger there shows the freeze state).}

\paragraph{Is the beta-min reasonable?} The beta-min \lab{BM}, at the usual order $\sigs\rn$, has the same form as \SL's proof rates, with a fixed multiple of
$\log n$ where theirs has a diverging multiple of $\bar r_n\log P_n$ (supp.\ p.~15). \SL's Theorem~1 bounds the
$\ell_1$ norm of the coefficients through (3) (p.~950), and (8) bounds the part outside the best $r_n$ coefficients
(p.~952). Without any lower bound, specificity can fail
(Remark~\ref{rem:no-betamin}).
\emph{Precedents:} \citet[Theorem~5 and Corollary~1, pp.~13--14]{CastilloSchmidtHieberVanDerVaart2015} and
\citet[Theorem~1, display (8), p.~9]{YangWainwrightJordan2015} require the nonzero coefficients to exceed a multiple of
$\sqrt{\log p/n}$ (in the first, with column norms of order $\sqrt n$, compatibility numbers bounded below and a bound on
the true model size through the mutual coherence; in the second, for the coordinates it recovers); when $p\ge n^a$ for some fixed $a>0$ our order $\sigs\rn$ is the same, and when $\log p=o(\log n)$ ours is the larger requirement.

\subsubsection{Further posterior results}

\begin{remark}[The posterior miss at moderate $n$; a heuristic, not a theorem]\label{rem:moderate-n}
The posterior miss is carried mainly by the models
$\Sz\cup\{j\}$ that add one null predictor; each is light, but together, since the fit factor $e^{T_j^2/2}$ ($T_j^2$
approximately $\chi^2_1$) has no finite mean, they carry about $(\sz+1)\sqrt{\log p/(\pi n)}$ times the weight of $\Sz$:
the expression of Theorem~\ref{rem:sharp}(b), here as a ratio to the weight of $\Sz$. With $\sz=5$, $n=200$ and $p=10^4$ this is about $0.7$. In single
synthetic draws (AR(0.6), $\sz=5$) the measured total weight of these models, relative to that of $\Sz$, was $1.0$--$1.2$
at $n=200$ and $p=10^4$--$10^6$, and $0.52$ at $n=400$, $p=10^4$. So at simulated sizes the posterior is spread over $\Sz$
and many one-predictor extensions of it.
\end{remark}

\why{Without the beta-min condition, the posterior still rejects over-fitted models and still keeps every true predictor
that is strong enough; specificity, however, can fail.}

\begin{remark}[Without beta-min]\label{rem:no-betamin}
Under the other assumptions of Theorem~\ref{thm:post-fixed}, the posterior mass of models that strictly contain $\Sz$ is
$\OP(\rn)$, and every true predictor with $\theta_j^{*2}\ge c\,\sigs^2\log n/(n\varphi_0^2)$, $c>c_0(\sz,A)$, has
$1-\PIP_j=\OP(\rn)$. Specificity can fail without beta-min: a numerical counterexample, reported with the proof, shows this.
\end{remark}

\status{Part of the posterior concentration result, with the status of the in-probability form of
Theorem~\ref{thm:post-fixed}.}

\why{It replaces the eigenvalue floor away from the true model, and the beta-min condition, by a floor on the true model
alone and a margin on the unexplained signal; it keeps the vanishing floor \as{NC}.}

\begin{localassumption}[Floor on the true model and an identifiability margin]\label{asm:margin}
\begin{asmitems}
  \item $\lmin(\Rhat_{\Sz})\ge\varphi_0^2$ for a fixed $\varphi_0\in(0,1]$.
  \item Every model $S\in\Om$ that misses a true predictor and has at most $j_0:=\lceil(2A+1)\sz\rceil+1$ predictors
    outside $\Sz$ leaves unexplained signal $\norm{(I-P_S)\mu}^2\ge c_{\mathrm{id}}\,\sigs^2\log n$, for a fixed
    $c_{\mathrm{id}}>c_1(\sz,A)$, an explicit constant (Appendix~\ref{app:constants}).
\end{asmitems}
\end{localassumption}

\begin{theorem}[Posterior concentration under an identifiability margin]\label{thm:post-margin}
Let the design be fixed and assume \as{GL}, \as{DIM}, \as{SP}, \as{BC}, \as{NC} and Assumption~\ref{asm:margin}. Then
the conclusions~\eqref{eq:post-op}--\eqref{eq:post-exp} of Theorem~\ref{thm:post-fixed} hold, the expectation bound with
constants depending only on $(\sz,A,\sigs^2,\bar\theta,\varphi_0,g_0,c_{\mathrm{id}})$.
\end{theorem}

\status{The in-probability conclusions: \textbf{proved and independently checked}. The expectation bound:
\textbf{proved and independently checked}, as the expectation bound of Theorem~\ref{thm:post-fixed} (its
lower-bound proposition checked on 2026-11-17). As there, the registry strings still read ``conditional on'' the same
checked lemma on random Gaussian designs. Registry: Appendix~\ref{app:status}, row 1.1.}

Under \as{EF} and \as{BM}(a), Assumption~\ref{asm:margin} holds with $c_{\mathrm{id}}=C_\beta^2\varphi_0^2(1-1/n)$. So
Theorem~\ref{thm:post-margin} contains Theorem~\ref{thm:post-fixed} when $C_\beta^2\varphi_0^2>c_1$; when
$c_0<C_\beta^2\varphi_0^2\le c_1$, Theorem~\ref{thm:post-fixed} applies and Theorem~\ref{thm:post-margin} is not known to
apply: the two coexist, and neither contains the other.

\why{It says what the posterior selects when some true predictors are missing from the data: the predictors used by the
best linear predictor built from the available ones.}

In this result the method is run on the response $\yr$ and a set $\mathcal C$ of $p$ candidate predictors, with rows
$x_{\mathcal C,i}$ collected in $x_{\mathcal C}$; $\Sigma$ and $R$ are the covariance and the correlation matrix of a
candidate row.

\begin{localassumption}[Candidate-only fitting]\label{asm:omitted}
\begin{asmitems}
  \item \emph{Model.} The candidate rows are i.i.d.\ $\Normal_p(m_{\mathcal C},\Sigma)$, and, given them, the responses are
    independent with $\yr_i\mid x_{\mathcal C}\sim\Normal(\alpha'+x_{\mathcal C,i}^\top\beta^{\mathcal C},\sigma^2_{\mathcal C})$,
    $\sigma^2_{\mathcal C}>0$. We write $\Sz^{\mathcal C}:=\supp\beta^{\mathcal C}$ and $s_{\mathcal C}:=\abs{\Sz^{\mathcal C}}$.
  \item \emph{Floor.} $\lmin(R_T)\ge c_{\mathcal C}$ for every $T\subseteq\mathcal C$ with
    $\abs{T}\le3(A+1)(s_{\mathcal C}+1)\log n$, for a fixed $c_{\mathcal C}>0$.
  \item \emph{Signal relative to the residual variance.} $\beta^{\mathcal C\top}\Sigma\beta^{\mathcal C}\le
    M_{\mathcal C}\sigma^2_{\mathcal C}$ for a fixed $M_{\mathcal C}$, and $n\min_{j\in\Sz^{\mathcal C}}\Sigma_{jj}(\beta^{\mathcal
    C}_j)^2/(\sigma^2_{\mathcal C}\log n)\to\infty$.
\end{asmitems}
\end{localassumption}

\begin{theorem}[Omitted variables]\label{thm:omitted}
Assume \as{DIM} (with $p$ the number of candidates), Assumption~\ref{asm:omitted}, and that $s_{\mathcal C}\ge1$ is fixed. The
quantities $m_{\mathcal C}$, $\Sigma$, $\beta^{\mathcal C}$ and $\sigma^2_{\mathcal C}$ may vary with $n$; no condition beyond
Assumption~\ref{asm:omitted} is placed on how they vary. Then
\[
  \E\big[1-\pi(\Sz^{\mathcal C}\mid y)\big]=O\Big(\sqrt{\frac{\log n}{n}}\Big),\qquad
  \P\big(\HPM\neq\Sz^{\mathcal C}\big)=O\Big(\sqrt{\frac{\log n}{n}}\Big).
\]
\end{theorem}

\status{\textbf{Proved and independently checked}, using the posterior concentration result and two lemmas on Gaussian
designs that are independently verified; since 2026-11-17 the registry no longer lists the lower-bound proposition among
its conditions (that proposition is now checked). The registry strings still read ``conditional on'' inputs that are
themselves proved and checked or verified: the lemma on random Gaussian designs and, for one form, the two Gaussian-design
lemmas. Registry: Appendix~\ref{app:status}, row 1.4.}

\begin{example}[Omitted Gaussian predictors]\label{ex:omitted}
If candidates and omitted predictors are jointly Gaussian, the response given the candidates is Gaussian with a linear
mean, and $\beta^{\mathcal C}=\beta^*_{\mathcal C}+\Sigma_{\mathcal C\mathcal C}^{-1}\Sigma_{\mathcal C\mathcal M}\beta^*_{\mathcal M}$:
each candidate absorbs the part of the omitted signal it can predict, and the rest is added to the noise variance
$\sigma^2_{\mathcal C}$. This is an example of Assumption~\ref{asm:omitted}(a), not a claim that the assumption holds in
ordinary data.
\end{example}

\subsection{The distribution}\label{sec:supp-dist}

\why{It shows what the collinearity condition of Result~2 buys: dropping it costs one logarithmic factor in the radius,
and the posterior also puts almost no mass on large models.}

\begin{proposition}[Level-$M$ form, and the case without the collinearity condition]\label{prop:dist}
\textup{(i)} There is a constant $C_{\mathrm s}$, depending only on $(s,A,L,c_\psi,\bar b)$, such that the following holds. The
constants of Theorem~\ref{thm:dist} serve every level $M$ with $M_0\le M\le\sqrt n/\log n$, and $M$ may depend on $n$:
whenever $n\ge n_0$ and Assumption~\ref{asm:dist} holds with the window $q_n$ of this $M$,~\eqref{eq:dist-i} holds, and also
\begin{equation}\label{eq:dist-ii}
  \P^*\Big[\,\Pi\big(\abs{S}>s+C_{\mathrm s}M\ \big|\ \Dn\big)>n^{-c_{\mathrm H} M}\Big]\ \le\ n^{-c_{\mathrm H} M}.
\end{equation}
The two exceptional sets can be taken to be the same, and the bounds are uniform over all laws of the data that satisfy
Assumption~\ref{asm:dist} with the given constants.
\textup{(ii)} If Assumption~\ref{asm:dist}\ref{asm:dist-nec} is dropped and Assumption~\ref{asm:dist}\ref{asm:dist-subg} is
required only for single predictors, then~\eqref{eq:dist-i} holds with radius $\sqrt{M\log n\cdot\log(\praw\vee n)/n}$,
and~\eqref{eq:dist-ii} with $s+C_{\mathrm s}M\log n$ in place of $s+C_{\mathrm s}M$.
\end{proposition}

\status{\textbf{Proved and independently checked} (part of the same result as Theorem~\ref{thm:dist}). Registry:
Appendix~\ref{app:status}, row 2.1.}

\why{With a fixed eigenvalue floor and every true signal above the detection scale, the distribution is learned at the
parametric rate.}

\begin{theorem}[Parametric rate when the true predictors are detectable]\label{thm:dist-param}
Assume Assumption~\ref{asm:dist} with the window $q_n$ of level $M=1$, with exactly $s\ge1$ nonzero coordinates of $\bstar$;
that $\lmin(\Corr(x_T))\ge\psi$ for a fixed $\psi>0$ and every set $T$ of at most $(2A+3)s+2$ predictors; and that
$\min\{\abs{\beta^*_j}\sd(x_j)/\sigs:\beta^*_j\neq0\}\ge b_n$, where $b_n>0$ and $b_n^2n/\log n\to\infty$. Then
$\Pi(d_H(f,f^*)>M_n/\sqrt n\mid\Dn)\to0$ in $\P^*$-probability, for every sequence $M_n\to\infty$.
\end{theorem}

\status{\textbf{Proved and independently checked}. Registry: Appendix~\ref{app:status}, row 2.1.}

\why{The distribution guarantee carries over to the list of models that a completed run returns, provided that list
holds a non-negligible share of the posterior, a premise that is not proved here; it is therefore stated as a conditional
remark.}

A completed run returns a list $\mathcal T$ of at most $40p+1$ distinct models with their log scores; let
$\pi(\mathcal T\mid\Dn)$ be the posterior probability of the list, and $\hat\Pi$ the posterior restricted to it (the models
of $\mathcal T$ weighted in proportion to their exponentiated log scores, each with its exact conditional posterior of
$(\beta,\sigma^2)$; $\hat\Pi$ is not computed by the method).

\begin{remark}[The output of the sampler, conditionally on the mass of its list]\label{cor:dist-output}
Assume the conditions of Proposition~\ref{prop:dist}(i), and that the returned log scores are computed in exact
arithmetic. With $\P^*$-probability at least $1-n^{-c_{\mathrm H}M}$, every completed run, whatever its random numbers, with
$\pi(\mathcal T\mid\Dn)\ge n^{-c_{\mathrm H}M/2}$ satisfies $\hat\Pi(d_H(f,f^*)>\sqrt{M\log n/n})\le n^{-c_{\mathrm H}M/2}$. If a
completed run scored $\Sz$ and $\pi(\Sz\mid\Dn)>\tfrac12$, then the list contains $\Sz$, $\Shat=\Sz$, and $\pi(\mathcal
T\mid\Dn)>\tfrac12$.
\end{remark}

\status{\textbf{Proved and checked, conditional on} exact arithmetic only; the two results on the coded sampler it also
used were line-checked and recorded on 2026-11-17, after the freeze, so they are no longer conditions. The description of
the output was checked against the code. Registry: Appendix~\ref{app:status}, row 2.1.}

Whether the list carries enough posterior probability is a property of the run; this remark does not assert it, and
in general nothing in this paper proves it. On the designs of Theorem~\ref{thm:reported} the list holds almost all the
posterior mass (Corollary~\ref{cor:list}); that this yields the premise here is not claimed.

\subsection{The method}\label{sec:supp-method}

\why{It explains why the sampler starts in the right place: without leakage, the start law already concentrates on the
true model.}

\begin{theorem}[The start law concentrates]\label{thm:start}
Assume \as{GR}, \as{DIM} with $p>\sz$, \as{FSt}, \as{V}, \as{NL}, and $\lmin(\Sigma_{n,S})\ge\varphi_{\mathrm b}^2$ for every
$S$ with $\abs{S}\le16+\sz$, for a fixed $\varphi_{\mathrm b}>0$. Then $\P(\pi_0(\Sz)\ge1-n^{-c})\to1$ for every
$c<\tfrac12$. If moreover $\Cov(W_l,W_j)=0$ for every null $l$ and true $j$, then $1-\pi_0(\Sz)=\OP(\rn)$. The exponent
$\tfrac12$ cannot be improved: $1-\pi_0(\Sz)\ge\tfrac13(n+1)^{-1/2}$ for every valid data set with $n\ge4$ and $\sz<K$.
\end{theorem}

\status{\textbf{Proved and independently checked}. Registry: Appendix~\ref{app:status}, row 3.3.}

With leakage, the start law concentrates instead on the models between the true model and the leakage set: keep the
other assumptions of Theorem~\ref{thm:start}, and let \as{NL} hold with $\Sz$ replaced by a deterministic $Q\supseteq\Sz$ with $\abs{Q}\le K$, so
that only predictors outside $Q$ are required not to leak. Then $\P\big(\pi_0(\Lat(Q))\ge1-(\abs{Q}+1)n^{-1/2}e^{C\sqrt{\log n}}\big)\to1$.

\why{It is the only reported-model statement that does not assume Gaussian rows: its conditions are on the sample
$(X,\tstar)$.}

For a set $U$ and a predictor $l\notin U$, the \emph{signal gain} $\delta^2_{U,l}:=\norm{(P_{U\cup\{l\}}-P_U)\mu}^2$ is the
increase in explained mean signal when $x_l$ joins $X_U$; the \emph{residual scale} is $\tau^2_U:=\sigs^2+\norm{(I-P_U)\mu}^2/(n-1)$;
and $\hat\beta_{U,l}:=(X_U^\top X_U)^{-1}X_U^\top x_l$, padded with zeros.

\begin{localassumption}[Fixed-design visibility and leakage]\label{asm:reported-fixed}
\begin{asmitems}
  \item\label{asm:rf-floor} \emph{Floors.} $\lmin(\Rhat_T)\ge\varphi_0^2$ for every $T$ with $\abs{T}\le16$ and for every
    $T=Q\cup\{l\}$, $l\in\set{p}$, with $Q$ as in~\ref{asm:rf-leak}.
  \item\label{asm:rf-vis} \emph{Visibility.} There is a deterministic $\omega_n\to\infty$ such that
    $\delta^2_{U,j}\ge\omega_n\log n$ for every $j\in\Sz$ and every $U\not\ni j$ with $\abs{U}\le15$.
  \item\label{asm:rf-leak} \emph{Bounded leakage.} There are a deterministic $Q\supseteq\Sz$ (with $\abs{Q}\le K$ and
    ($Q=\set{p}$ or $K\ge2\abs{Q}+32\sz+1$) for large $n$) and a fixed $\eta\in[0,\eta_A)$,
    $\eta_A:=(\sqrt{2A+1}-\sqrt{2A})^2$, such that $\delta^2_{U,l}\le\eta\,\tau^2_U\log n$ for every $l\notin Q$ and every
    $U\not\ni l$ with $\abs{U}\le15$.
  \item\label{asm:rf-size} \emph{Size of the leakage set.} One of: (d1) $\abs{Q\setminus\Sz}\le(1-\delta_Q)\log n$ for a
    fixed $\delta_Q\in(0,1)$; or (d2) $\abs{Q}\le n^b$ with $b<\tfrac12-\tfrac\eta2-\sqrt{2A\eta}$,
    $\theta^{*2}_{\min}n/(\abs{Q}+\log n)\to\infty$, and local regressions: for every $l$ a deterministic set
    $N_l\subseteq Q\setminus\{l\}$ with $\abs{N_l}\le h\log_2n$, for a fixed $h$ with $h+b<\tfrac12$, and a deterministic $L_n\to0$ with
    $L_n^2\log(\abs{Q}+1)=o(\log n)$, such that $\norm{\hat\beta_{T,l}-\hat\beta_{T\cap N_l,l}}_1\le L_n$ for every
    $T\in\Lat(Q)$ with $\abs{T}<K$ and $l\notin T$.
\end{asmitems}
\end{localassumption}

\begin{theorem}[The reported model equals the true model: fixed design]\label{thm:reported-fixed}
Let the design be fixed, with no raw predictor constant in the sample (an input condition of the delivery result; then
$p=\praw$). Assume \as{GL}, \as{DIM}, \as{SP}, \as{BC},
\as{EF}, \as{NC}, \as{BM}(b), Assumption~\ref{asm:reported-fixed} and Convention~\as{C1}. Then the conclusions of
Theorem~\ref{thm:reported} hold: there are data events $\mathcal G_n$ with $\P(\mathcal G_n)\to1$ on which every
admissible run that produces output reports $\Shat=\Sz$, for every seed and calibration outcome; and the delivery
clauses hold under \as{C2} on $\Dom$ and under \as{C2p} for every admissible $n\ge3$.
\end{theorem}

\status{\textbf{Proved and independently checked}, with no open condition: the version proved in the registry is a stronger form of
this statement. Registry: Appendix~\ref{app:status}, row 3.2.}

\why{It states in one place when Theorem~\ref{thm:reported} applies to block and autoregressive designs.}

\begin{corollary}[Block and autoregressive designs]\label{cor:famc}
In the setting of Theorem~\ref{thm:reported}, suppose that for all large $n$ the predictors are correlated either in
independent blocks of bounded size (uniformly well conditioned) or autoregressively with a fixed $\rho\in(-1,1)$, and that
a fixed number of true coefficients have absolute values bounded away from $0$ and $\infty$. Suppose also that the true
predictors cannot cancel one another: (a) one true predictor per block; or (b) signs aligned with the correlation (for AR,
equal signs if $\rho>0$ and signs alternating with the column index if $\rho<0$; for a block with several true predictors,
after flipping some columns' signs all its partial correlations are nonnegative and the flipped true coefficients share
one sign); or (c) for AR only, the true predictors are spread apart by an explicit separation bound; or (d) the coefficient
vector stays a fixed distance from a closed set $B$ of cancelling directions ($B$ is null for AR designs, and possibly not
null for general block designs). Then, on data events of probability tending to one, every admissible run that returns a
result reports exactly the true model, whatever its seed. Under the delivery conventions of Theorem~\ref{thm:reported},
every run returns a result with probability one, so the true model is recovered with probability tending to one, with
that theorem's qualifications (lifted integer limits, bounded domain, time guard).
\end{corollary}

Not covered: coefficients that cancel, such as two adjacent autoregressive true predictors with coefficients $1$ and $-2$
at correlation $0.5$. Our designs meet the conditions: block designs with one true predictor per block, and the AR($0.5$)
and AR($0.6$) layouts, whose true predictors are spread apart. The constants degrade as $\abs\rho\to1$.

\status{\textbf{Proved and independently checked} (recorded on main at 8b4d3101): an assembly of recorded results, with no
new assumption. Registry: claim W3-FAMC:Cor-FAMC.}


\subsection{Estimates from the returned list (Result 3)}\label{sec:supp-list}
 For a nonempty set $\mathcal T\subseteq\Om$ of models, the
\emph{list-renormalised} quantities are the posterior renormalised to $\mathcal T$ and its total weight:
\[
  Z_{\mathcal T}:=\sum_{S\in\mathcal T}e^{F(S)},\qquad
  \mathrm{pip}_{\mathcal T}(j):=\frac{1}{Z_{\mathcal T}}\sum_{S\in\mathcal T,\,S\ni j}e^{F(S)},\qquad
  \mathrm{mpm}_{\mathcal T}:=\{j:\mathrm{pip}_{\mathcal T}(j)>\tfrac12\}.
\]
They are functions of the list $\mathcal T$ that a run returns and of its log scores, not outputs of the method as
Section~\ref{sec:outputs-method} describes it. Below, $\varepsilon_n$ is the posterior miss (not the noise).

\begin{corollary}[List-renormalised estimates on the designs of Theorem~\ref{thm:reported}]\label{cor:list}
Under the assumptions of Theorem~\ref{thm:reported}, there are data events of probability tending to one on which every
admissible run that produces output, for every seed and calibration outcome, reports $\Shat=\Sz$ and has $\Sz$ in its
returned list $\mathcal T$ (on exact enumeration, in the enumeration's list where it is returned); and then, with
$\varepsilon_n:=1-\pis{\Sz}$,
\begin{gather}
  \max_j\big|\mathrm{pip}_{\mathcal T}(j)-\ind{j\in\Sz}\big|\le\varepsilon_n,\qquad
  \max_j\big|\mathrm{pip}_{\mathcal T}(j)-\PIP_j\big|\le\varepsilon_n,\label{eq:list-pip}\\
  1-\varepsilon_n\ \le\ Z_{\mathcal T}/Z\ \le\ 1,\label{eq:list-z}\\
  \mathrm{mpm}_{\mathcal T}=\Sz\quad\text{when }\varepsilon_n<\tfrac12 .\label{eq:list-mpm}
\end{gather}
Since $\varepsilon_n=\OP(\rn)$, these errors are $\OP(\rn)$ for runs that produce output; under Convention~\as{C2p} every
run does.
\end{corollary}

\status{\textbf{Proved and independently checked}, recorded on 2026-11-18, after the freeze, with no condition: the code
reading of the method commit that adds the list outputs was done (it is pinned at an earlier core commit; its re-pin to the
final method tag is a later step). Its inputs, the masters of Theorem~\ref{thm:reported} and the posterior concentration
result, are proved and independently checked.}

If the list contains the true model and the true model carries almost all the posterior mass, then renormalising
to the list loses almost nothing. The pooled estimates $\PIPhat_j$, $\MPMhat$ and $\Zhat$ of~\eqref{eq:pool} are not
covered by it; their consistency is proved only on designs without leakage (Remark~\ref{thm:outputs}).

On the part of its event where $\pis{\Sz}\ge n^{-c_{\mathrm H}M/2}$, Corollary~\ref{cor:list} supplies, pathwise, the
list-mass premise of Remark~\ref{cor:dist-output} (the list then holds at least that mass); the probability form under that
remark's law needs a raw-scale bridge that is open.

\subsection{The reported estimates are consistent, on designs without leakage}\label{sec:outputs}

\why{It is the only result on the consistency of the reported inclusion probabilities, median model and evidence
estimate; it needs the condition \as{NL}, which, of our designs, only the independent ones satisfy, so it is stated
as a remark.}


\begin{assumption}{NL}{No leakage, as seen by the start law}
Every null predictor $l\notin\Sz$ has a deterministic \emph{neighbourhood} $\Hl_l\supseteq\{l\}\cup\Sz$ with
$\abs{\Hl_l}\le e^{\kappa_H\sqrt{\log n}}$ for a fixed $\kappa_H$, such that, for every $U\not\ni l$ with $\abs{U}\le15$ and
$U_0:=U\cap\Hl_l$: (i) \emph{no leak:} $\abs{\Cov(W_l,W_j\mid W_{U_0})}\le C_Ln^{-1/2}$ for every $j\in\Sz\setminus U_0$,
with a fixed $C_L$; and (ii) \emph{locality:} $\abs{\Cov(W_l,W_v\mid W_{U_0})}\le r^\circ_n$ and
$\abs{\Cov(W_j,W_v\mid W_{U_0})}\le r^\circ_n$ for every $v\in U\setminus\Hl_l$ and $j\in\Sz\setminus U_0$, where
$r^\circ_n(n\log n)^{1/4}\to0$. \hfill [\SA{pe}]
\end{assumption}

The remark concerns runs that use the sampler; enumeration is not covered. On the event $\Acal$ of
Section~\ref{sec:conventions}, $\Zhat$, $\PIPhat$ and
$\MPMhat$ denote the pooled values the four islands would produce if the run did not stop; on completed runs they
are the reported values. Let $\rho_n:=n^{-1/2}e^{C_\rho\sqrt{\log n}}$, with $C_\rho$ depending only on the fixed quantities of the
assumptions.

\begin{remark}[The reported estimates are consistent]\label{thm:outputs}
Assume \as{GR}, \as{DIM} with $p>\sz$, \as{FSt}, \as{EFS}, \as{V}, \as{NL}, and Conventions~\as{C1} and~\as{C3}. Then,
under the joint law of the design, the noise and the method's random numbers,
\begin{gather}
  \ind{\Acal}\Big(\big|\Zhat/Z-1\big|+\big|\log\Zhat-\log Z\big|+\max_{j\in\set{p}}\big|\PIPhat_j-\ind{j\in\Sz}\big|\Big)
    =\OP(\rho_n),\label{eq:outputs-rate}\\
  \P\big(\Acal\cap\{\MPMhat\neq\Sz\}\big)\to0,\qquad \P\big(\text{the run completes and }\Shat\neq\Sz\big)\to0,
    \label{eq:outputs-sel}
\end{gather}
and~\eqref{eq:outputs-rate} also holds for the reported values multiplied by the indicator of a completed run.
\end{remark}

\asmlist{(1)~\as{GR}, \as{FSt}, \as{EFS}, \as{V} as in Result~3; (2)~\as{DIM} [\SA{all}]; (3)~\as{NL} no leakage
[\SA{pe}]; (4)~Conventions \as{C1}, \as{C3} (Appendix~\ref{app:conventions}).}

\status{\textbf{Proved and independently checked}; the two results on the coded sampler it rests on were
line-checked and recorded on 2026-11-17, after the freeze, so the registry lists no condition. Registry:
Appendix~\ref{app:status}, row 3.6 (the ledger shows the freeze state, with the condition).}

\par\noindent\textit{What is not claimed.}\phantomsection\label{rem:outputs}
(i) No lower bound on $\P(\Acal)$ is asserted: the statements are about intersections with $\Acal$ or with completion, not
conditional on them.
(ii) The rate $\rho_n$ is the start law's rate; it is slower than the posterior rate $\rn$, by the factor
$e^{C_\rho\sqrt{\log n}}/\sqrt{\log n}$.
(iii) Pooling cannot make things worse: the pooled inclusion probabilities are never farther from the exact ones, in the
largest coordinate, than those of the worst island (proved and independently checked; unconditional since 2026-11-17, as
Remark~\ref{thm:outputs}; registry row 4.1).

\paragraph{Is the no-leakage condition reasonable?}\phantomsection\label{sec:a-pe}

\ahead{No leakage}{\lab{NL}}{In no regime do \SL{} have a condition of this kind: none of their theorems concerns a
start law. Only our independent designs satisfy it.}
\ev{With no conditioning, \lab{NL} asks $\abs{\Cov(W_l,W_j)}\le C_Ln^{-1/2}$ for every null predictor $l$ and true
predictor $j$. In AR(1) with correlation $0.5$ the neighbour of a true predictor has covariance $0.5$; in a block design
every block-mate of a true predictor has covariance at least $0.4$. For such designs the start law concentrates on the
models between the true model and the leakage set, not on the true model alone. On AR and BLK the present proof route
cannot be extended as $n\to\infty$, because the start law's mass on the true model then decays faster than any power of
$n$; at the sizes of the $10^6$ plan this asymptotic decay says nothing, and we have no bound on this mass there in either
direction. For S3, S4 and S6, a theorem under a weaker, group-relative form of \lab{NL} is a candidate; it is not
proved.}
\ro{Any null predictor correlated with a true one at a fixed level.}
\ds{\textbf{Inside:} the independent designs only. \textbf{Outside:} all locally correlated and factor designs, including AR, the
primary design of the $10^6$ plan.}

\why{It is the only recorded route to the reported inclusion probabilities and median model that does not need \as{NL};
its global-overlap premise is not yet derived from conditions on the design.}

\begin{remark}[An alternative route for the estimates]\label{rem:alt-route}
An argument that does not go through the evidence estimate gives the conclusions of Remark~\ref{thm:outputs} for the
inclusion probabilities and the median and highest posterior models under a global-overlap condition (stated in full with
the proof).
\end{remark}

\status{\textbf{Proved and independently checked}; its global-overlap premise is a named open premise. Registry:
Appendix~\ref{app:status}, row 3.6.}


\subsection{The calibration and the Markov moves in full}\label{sec:supp-cal}

\why{Every result holds for every outcome of the calibration, so the main text uses only three facts about it; this is the
procedure itself, for reproducibility.}

At tier $i$ the calibration's islands have $\Nc_i:=4^i\Nc_0$ particles, and the \emph{step ceiling} is
$\Tmax_i:=4^i\max(1000,\ell_n^2)$.

\paragraph{Stores of models.} Throughout the run the method keeps three stores of the models scored so far: the $40p+1$
best by $F$; the $20p$ best by $\gamma_0$; and the $20p$ best by $\gamma_1=e^{F}/m(\emptyset)$, that is, by $F$. The two
stores of $20p$ feed the candidate set below. \emph{The best model found so far} is the model with the highest $F$ among
the models recorded so far: during the calibration, the block-table models (recorded before the pilot starts) and those
scored in the calibration; in the final run, the block-table models, those scored in the final run, and the models passed
on by the calibration.

\begin{center}
\fbox{\begin{minipage}{0.93\linewidth}\small
\textbf{Algorithm 1: calibration.}
\begin{enumerate}[leftmargin=1.6em,itemsep=1pt,label=\arabic*.]
  \item Set the tier $i\leftarrow0$.
  \item \emph{Pilot.} Run one island of $\Nc_i$ particles (reweight, resample, move, as in Section~\ref{sec:algorithm})
    with an adaptive schedule. From the current $\xi$, the next temperature is $\xi'=1$ if that step keeps an effective
    sample size of at least $0.90\,\Nc_i$; otherwise it is found by 60 bisections on $(\xi,1]$. If the step is shorter than
    $1/\Tmax_i$, take $\xi'=\min(1,\xi+1/\Tmax_i)$ instead and count a \emph{forced step}. At $\xi=1$, run $\Lbi$ burn-in rounds and
    four measured rounds, then a local search (move cap $\ell_n$) from the best model found so far.
  \item \emph{Grid.} Build a grid from the pilot's temperatures on the candidate set $\mathcal V$ (below). If the
    construction fails, record a trigger and skip step~4.
  \item \emph{Rebuild loop.} Repeat at most $\ell_n$ times: run four diagnostic islands of $\Nc_i$ particles on the
    current grid (each as in Section~\ref{sec:algorithm}, with $\Lbi$ burn-in and four measured rounds); run a local search
    (move cap $\ell_n$) from the best model found so far; recompute $\mathcal V$, and rebuild the grid from the pilot's
    temperatures. Stop when the rebuilt grid is identical to the previous one (the grid has \emph{converged}), or when a
    rebuild fails.
  \item \emph{Triggers.} If $i=0$ and any trigger fires (list below): repeat the check before the run, set $i\leftarrow1$,
    and go to step~2. The scored models, the three stores and the score cache carry over; the tier-$0$ temperatures do
    not, and tier~1 uses its own random streams. At tier~1, triggers are only recorded.
  \item \emph{Fallback.} If the grid construction failed at tier $1$, use the last grid built successfully at tier $1$,
    or else the tier-$1$ pilot's own temperatures. (An unconverged grid is used as it is.)
  \item \emph{Freeze} the grid, the tier $i$, the number of particles $N_i$ and $\Lbi$, and pass the three stores and the
    calibration's best model to the final run, as candidates for the reported model (they are not particles). With the
    fallback on (the default), a grid always exists at this point; if it does not, stop the run without estimates.
\end{enumerate}
\end{minipage}}
\end{center}

\paragraph{The candidate set and the grid.} The candidate set $\mathcal V$ is the empty model together with the two stores
of $20p$ models (the best by $\gamma_0$ and the best by $F$). On $\mathcal V$, let
$\psi(u):=\log\sum_{S\in\mathcal V}\gamma_0(S)e^{u\,\Delta(S)}$. A step from temperature $u$ to $v$ has \emph{restricted effective
sample size}
\[
  \rho(u,v):=\exp\{2\psi(v)-\psi(u)-\psi(2v-u)\},
\]
the effective-sample-size fraction that the step would leave if the particles were the models of $\mathcal V$ with
weights proportional to $\gamma_u$. The grid \emph{refines} the pilot's temperatures. Starting at $u=0$, let $b$ be the
smallest pilot temperature above $u$. If $\rho(u,b)\ge0.90$, the next grid point is $b$; otherwise it is the largest
point of $(u,b]$ that passes, found by 60 bisections. The construction continues from the new point, again aiming at the
smallest pilot temperature above it, until it reaches $1$; so every pilot temperature is in the grid, and points are
inserted where a pilot step fails the check on $\mathcal V$. The construction \emph{fails} if a step
makes no numerical progress, if an appended temperature does not pass the $0.90$ check (temperatures within $10^{-12}$ of
$1$ are set to $1$), or if more than $\Tmax_i$ steps are needed.

\paragraph{Triggers.} At tier $0$ the calibration moves to tier $1$ if any of these occurs: \emph{island disagreement}
(E1: some diagnostic island's best log score is more than $1$ below the best score found so far); \emph{evidence
spread} (E2: the sample standard deviation of the four islands' log-evidence estimates exceeds $1$); \emph{late
improvement} (E3: a calibration local search, the pilot's included, ended with a strictly higher best log score than the
best before it started); a forced pilot step; a failed grid construction; a grid
that did not converge; or a local search that was still improving when it reached its move cap. E1 and E2 are taken from
the last completed rebuild; E3 and the move-cap trigger count over the pilot and every rebuild.

\paragraph{When the calibration succeeds.} The calibration \emph{succeeds} when it freezes a grid, built by the main
construction or by the fallback. Without a trigger, tier $0$ freezes its own grid. The calibration's own estimates are
archived and never reported.


\paragraph{The Markov moves.} One \emph{round} applies $p+B+2$ moves to a particle, in an order drawn uniformly at random
afresh for each particle and each round; each move leaves $\pi_\xi$ invariant.
\begin{itemize}[leftmargin=1.5em,itemsep=1pt]
  \item \emph{Inclusion update of predictor $j$} (one for each $j\in\set{p}$; a Gibbs update): with
    $a:=\gamma_\xi(S\cup\{j\})/\gamma_\xi(S\setminus\{j\})$, include $j$ with probability $o_j:=a/(1+a)$, or $o_j:=0$ if
    including $j$ would exceed the cap $K$. The probability $o_j$, not the drawn indicator, is what the inclusion-probability
    estimate uses (step~3 below).
  \item \emph{Uniform swap:} pick a member $j$ of the model and a non-member $l$ uniformly at random; propose
    $S':=(S\setminus\{j\})\cup\{l\}$, and accept with probability $\min\{1,\gamma_\xi(S')/\gamma_\xi(S)\}$. (If
    the model is empty or equal to $\set{p}$, the move does nothing.)
  \item \emph{Neighbour swaps} ($B$ of them): pick a member $j$ uniformly, and a non-member $l$ uniformly among $j$'s
    neighbours in the neighbour graph; propose the swap, and accept with the Metropolis--Hastings probability, which
    includes the correction $\abs{\mathcal N(j)\setminus S}/\abs{\mathcal N(l)\setminus S'}$ for the different
    neighbourhood sizes. (If the model is empty, or $j$ has no neighbour outside it, the move does nothing.)
  \item \emph{Independence move:} propose a fresh draw $S'$ from $\pi_0$, and accept with probability
    $\min\{1,\exp\{\xi\,(\Delta(S')-\Delta(S))\}\}$.
\end{itemize}
Swaps keep the model size; the inclusion updates and the independence move can change it. Every move stays in $\Om$:
the inclusion updates enforce the cap, and $\pi_0$ lives on $\Om$.



\subsection{Remarks moved from the main text}\label{sec:supp-moved}

\begin{remark}[Smaller sets, in expectation]\label{rem:thm2x}
Take i.i.d.\ Gaussian rows (with any positive variances), and ask the floor only on sets of at most
$3(A+1)(\sz+1)\log n$ predictors, on the correlation matrix of a row. Let the noise variance $\sigma^2$ depend on $n$, ask
the signal energy $\bstar{}^\top\Sig\bstar$ to be at most a fixed multiple of $\sigma^2$, and ask
$n\min_{j\in\Sz}\Sigma_{n,jj}\beta^{*2}_j/(\sigma^2\log n)\to\infty$ ($\Sig$ and $\bstar$ may also vary with $n$). Then $\E[1-\pis{\Sz}]=O(\rn)$, and the HPM and MPM error probabilities and
the expected maximal PIP error are $O(\rn)$. This is Theorem~\ref{thm:omitted} with no omitted predictors, with
Lemma~\ref{lem:summaries}. Against \as{EFS} as stated, it asks the floor on smaller sets. The result behind the Gaussian
case of Theorem~\ref{thm:post-random} can let the floor on its larger sets vanish slowly; this one cannot.
\end{remark}

\status{\textbf{Proved and checked, conditional} on three inputs of Theorem~\ref{thm:omitted}: one result of the posterior
concentration family and two earlier-wave lemmas on Gaussian designs. Recorded on 2026-11-17, after the freeze.}

\begin{remark}[When visibility holds]\label{rem:vis}
Visibility holds with a margin that does not depend on $n$ (i) with \emph{aligned signs}: when the precision matrix,
after flipping the signs of some columns, has nonpositive off-diagonal entries and, within each block, the true
coefficients' signs agree with those flips (independent columns; autoregressive correlation (moderate or not) or one factor with aligned
coefficients; blocks with the true predictors in different blocks); and (ii) in \emph{generic position}: for designs with
countably many local correlation patterns (finitely many block, factor or tile types; stationary designs with decaying
correlation, such as autoregressive, banded and Toeplitz), for every vector of true coefficients outside one fixed closed,
Lebesgue-null set of hyperplanes, the same for every $n$. Not covered: coefficient vectors on that set, such as strong-factor designs with mixed-sign coefficients, or adjacent
autoregressive true predictors with coefficients $1$ and $-2$ at correlation $0.5$. A design modelled on the riboflavin
gene-expression data is covered when its four true predictors are spread apart: they then lie in four different tiles
once $p\ge16{,}352$, with a margin of at least $0.1\min_k\abs{\theta^*_k}$. It is not covered when they are clustered
in one tile, where no usable margin can be certified.
\end{remark}

\status{Lemma VIS: \textbf{proved and independently checked}, recorded on 2026-11-17, after the freeze. It gives
sufficient conditions for \as{V}; it adds no assumption.}

\begin{remark}[On the assumptions]\label{rem:dist-asm}
No beta-min condition and no fixed lower bound on eigenvalues are assumed; jointly Gaussian predictors satisfy
Assumption~\ref{asm:dist}\ref{asm:dist-subg} with $L=1$. Without Assumption~\ref{asm:dist}\ref{asm:dist-nec} the radius grows
by one factor $\sqrt{\log(\praw\vee n)}$ (Proposition~\ref{prop:dist}). Display~\eqref{eq:dist-i}, with
Assumption~\ref{asm:dist}\ref{asm:dist-nec}, carries over to the list a completed run returns when that list holds
posterior mass at least $n^{-c_{\mathrm H}M/2}$, a condition the run does not check
(Remark~\ref{cor:dist-output}).
\end{remark}
